\section{Livrable 4 et 5 -- Cartographie mémoire et sécurités}

\begin{UPSTIactivite}[][Cartographie mémoire M340]
    \UPSTIquestion{Identifier les \textbf{informations à échanger} entre
    l'automate M340 et chacun des acteurs suivants :
    \begin{itemize}
        \item le robot (commandes envoyées, comptes rendus reçus) ;
        \item l'IHM (consignes opérateur, informations de supervision) ;
        \item le convoyeur (commandes de déplacement, informations
        capteurs).
    \end{itemize}}

    \UPSTIquestion{Pour chacune de ces informations, proposer un
    \textbf{type de données} adapté (booléen, entier, mot, structure, etc.)
    et une taille cohérente avec le besoin.}

    \UPSTIquestion{En vous appuyant sur le document ressource EtherNet/IP
    (objets Assembly notamment), proposer une \textbf{cartographie mémoire}
    au sein de l'automate M340 organisant ces échanges (zones d'entrées,
    zones de sorties, découpage par équipement).}
\end{UPSTIactivite}

\begin{UPSTIactivite}[][Sécurités]
    \UPSTIquestion{Identifier les principaux \textbf{risques} liés à
    l'exploitation de cette cellule robotique (zone de travail du robot,
    accès aux cuves, présence humaine, capacité unitaire du guichet, etc.).}

    \UPSTIquestion{Pour chaque risque identifié, proposer une ou plusieurs
    \textbf{sécurités} à mettre en œuvre (matérielles et/ou logicielles) et
    préciser leur rôle dans les modes de fonctionnement définis au livrable
    précédent.}
\end{UPSTIactivite}
